\subsection{Cycles associated with a module}

In this number, we denote by A a Noetherian ring.

{}{}Let Z(A) be the free Z-module with basis the set of irreducible closed subsets of Spec(A); for every irreducible closed subset Y of Spec(A), we denote by [Y] the corresponding element of Z(A). The elements of Z(A) are sometimes called cycles.

Let M be a finitely generated A-module. For every prime ideal  of A which is a minimal element of , we have \(0 < \text{long}_{A_{p}}(M_{p}) < \infty\) (IV, § 2, no. 5, Corollary 2 to Proposition 7 and § 1, no. 4, theorem 2); we set

\[
z (\mathbf {M}) = \sum_ {\mathfrak {p}} \operatorname{long} _ {\mathrm{A} _ {\mathfrak {p}}} (\mathrm{M} _ {\mathfrak {p}}). [ \mathrm{V} (\mathfrak {p}) ],
\]

where  ranges over the finite set of minimal prime ideals of Supp(M).

Remark. --- For every \(\mathfrak{p} \in \operatorname{Spec}(\mathbf{A})\), one has \(z(\mathbf{A}/\mathfrak{p}) = [\mathbf{V}(\mathfrak{p})]\). More generally, let M be a finitely generated A-module, and let \((\mathbf{M}_i)_{0 \leqslant i \leqslant n}\) be a composition series of M such that, for \(0 \leqslant i \leqslant n - 1\), the module \(\mathbf{M}_i / \mathbf{M}_{i+1}\) is isomorphic to \(\mathbf{A}/\mathfrak{p}_i\), where \(\mathfrak{p}_i\) is a prime ideal of A (cf. IV, § 1, no. 4, theorem 1); then we have \(z(\mathbf{M}) = \sum_{i \in \mathbf{J}} [\mathbf{V}(\mathfrak{p}_i)]\), where  is the subset of I consisting of those \(i\) such that \(\mathfrak{p}_i\) is a minimal element of \(\{\mathfrak{p}_0, \ldots, \mathfrak{p}_{n-1}\}\) (IV, § 1, no. 4, theorem 2 and § 2, no. 5, Remark 1).

For every integer \(d\) , denote by \(Z_{\leq d}\) (resp. \(Z_d\) , respectively \(Z^{\geq d}\) , respectively \(Z^d\) ) the sub-\(\mathbf{Z}\)-module of \(Z(A)\) generated by the elements  where \([V(\mathfrak{p})]\) is a prime ideal of A such that \(\mathfrak{p}\) (resp. \(A\), respectively \(\dim(A/\mathfrak{p}) \leqslant d\), respectively \(\dim(A/\mathfrak{p}) = d\)\(\operatorname{ht}(\mathfrak{p}) \geqslant d\)\(\operatorname{ht}(\mathfrak{p}) = d\)). The elements of \(Z_{d}\) (resp. \(Z^{d}\)) are said to be the cycles of dimension \(d\) (respectively, of codimension \(d\)). We obviously have

\[
\mathbf {Z} _ {\leqslant d} = \mathbf {Z} _ {\leqslant d - 1} \oplus \mathbf {Z} _ {d}, \quad \mathbf {Z} ^ {\geqslant d} = \mathbf {Z} ^ {\geqslant d + 1} \oplus \mathbf {Z} ^ {d}.
\]

Let C be the set of classes of finitely generated A-modules (A, VIII, § 3, no. 5), and for each integer d, let \(C_{\leq d}\) (resp. \(C^{\geq d}\)) be the subset of C consisting of the classes of finitely generated A-modules of dimension \(\leq d\) (resp. whose support has codimension \(\geq d\) in Spec(A)).

Lemma 1. --- Let M be a finitely generated A-module and let d be an integer.

\begin{enumerate}
\def\labelenumi{\alph{enumi})}
\tightlist
\item
  For M to be of type \(C_{\leq d}\) it is necessary and sufficient that \(z(\mathbf{M}) \in \mathbf{Z}_{\leq d}\) ; the projection \(z_{d}(\mathbf{M})\) of \(z(\mathbf{M})\) onto \(Z_{d}\) parallel to \(Z_{\leq d-1}\) is then given by
\end{enumerate}

\[
z _ {d} (\mathrm{M}) = \sum _ {\dim (\mathrm{A} / \mathfrak {p}) = d} \operatorname{long} _ {\mathrm{A} _ {\mathfrak {p}}} (\mathrm{M} _ {\mathfrak {p}}). [ \mathrm{V} (\mathfrak {p}) ].
\]

\begin{enumerate}
\def\labelenumi{\alph{enumi})}
\setcounter{enumi}{1}
\tightlist
\item
  For M to be of type \(C^{\geq d}\) it is necessary and sufficient that \(z(\mathbf{M})\in\mathbf{Z}^{\geq d}\) ; the projection \(z^{d}(\mathbf{M})\) of \(z(\mathbf{M})\) onto \(\mathbf{Z}^{d}\) parallel to \(\mathbf{Z}^{\geq d+1}\) is then given by
\end{enumerate}

\[
z ^ {d} (\mathbf {M}) = \sum _ {\operatorname{ht} (\mathfrak {p}) = d} \operatorname{long} _ {\mathrm{A} _ {\mathfrak {p}}} (\mathrm{M} _ {\mathfrak {p}}). [ \mathrm{V} (\mathfrak {p}) ].
\]

For M to be of type \(C_{\leq d}\), i.e. of dimension \(\leq d\), it is necessary and sufficient that, for every minimal prime ideal  of , one has \(\leq d\), which means that \(z(M) \in Z_{\leq d}\). Suppose that , and let \(\leq d\) such that \(p \in Spec(A)\); then either \(p \notin Supp(M)\), and therefore \(M_p = 0\), or \(p \in Supp(M)\), and { is a minimal element of }; the coefficient of {} in \(z(M)\) is in both cases \(_{A_p}(M_p)\), whence a). Part b) is proved analogously; note that a finitely generated module M is of type \(C^{\geq d}\) if and only if one has  for every prime ideal \(M_{\mathfrak{p}} = 0\) of height \textless{}.

By prop. 9, c) and remark 3 of no. 4, the sets \(\mathcal{C}_{\leq d}\) and \(\mathcal{C}^{\geq d}\) are hereditary (A, VIII, § 10, no. 1, def. 1), and one may consider the Grothendieck groups  and  corresponding to them (loc. cit., no. 2); for every A-module M of finite type in \(\mathcal{C}_{\leq d}\) (resp. \(\mathcal{C}^{\geq d}\)), we denote by  (resp. \(\mathcal{C}_{\leq d}\)\(\mathcal{C}^{\geq d}\){) the associated element in }{}\(_{\leq d}\) (resp. {}{}\(^{\geq d}\)\(\mathcal{C}_{\leq d}\)\(\mathcal{C}^{\geq d}\)). By Lemma 1, the functions \(z_{d}\) and \(z^{d}\) are additive; hence there exist (loc. cit., Prop. 3) homomorphisms

\[
\zeta_ {d}: \mathrm{K} (\mathbb {C} _ {\leqslant d}) \rightarrow Z _ {d}, \quad \zeta^ {d}: \mathrm{K} (\mathbb {C} ^ {\geqslant d}) \rightarrow Z ^ {d}
\]

such that \(\zeta_d([M]_{\leq d}) = z_d(M)\) for every finitely generated A-module of type \(C_{\leq d}\) and \(\zeta^d ([N]^{\geq d}) = z^d (N)\) for every finitely generated A-module N of type \(C^{\geq d}\). Moreover, since \(C_{\leq d-1} \subset C_{\leq d}\) and \(C^{\geq d+1} \subset C^{\geq d}\), we have canonical homomorphisms

\[
i _ {d}: \mathrm{K} (\mathbb {C} _ {\leqslant d - 1}) \rightarrow \mathrm{K} (\mathbb {C} _ {\leqslant d}) \quad \text {\text{ and }} \quad i ^ {d}: \mathrm{K} (\mathbb {C} ^ {\geqslant d + 1}) \rightarrow \mathrm{K} (\mathbb {C} ^ {\geqslant d}).
\]

With these notations:

PROPOSITION 10. --- The sequences of Z-modules and homomorphisms

\[
\mathrm{K} \left(\mathbb {C} _ {\leqslant d - 1}\right) \xrightarrow {i _ {d}} \mathrm{K} \left(\mathbb {C} _ {\leqslant d}\right) \xrightarrow {\zeta_ {d}} \mathrm{Z} _ {d} \longrightarrow 0
\]

\[
\mathrm{K} \left(\mathbb {C} ^ {\geqslant d + 1}\right) \xrightarrow {i ^ {d}} \mathrm{K} \left(\mathbb {C} ^ {\geqslant d}\right) \xrightarrow {\zeta ^ {d}} \mathrm{Z} ^ {d} \longrightarrow 0
\]

are exact.

We have \(\zeta_d \circ i_d = 0\) . By Lemma 1, for every \(\mathfrak{p} \in \operatorname{Spec}(A)\) such that \(\dim(A/\mathfrak{p}) = d\) , one has \(\zeta_d([A/\mathfrak{p}]_{\leq d}) = z_d(A/\mathfrak{p}) = [V(\mathfrak{p})]\) , hence the homomorphism \(\zeta_d\) is surjective. By IV, § 1, no. 4, th. 1, \(K(C_{\leq d})\) is generated by the \([A/\mathfrak{p}]_{\leq d}\) , where \(\mathfrak{p} \in \operatorname{Spec}(A)\) and \(\dim(A/\mathfrak{p}) \leqslant d\) ; consequently, every element \(\xi\) of \(K(C_{\leq d})\) can be written \(\xi = i_d(\eta) + \sum_{i=1}^{k} n_i[A/\mathfrak{p}_i]_{\leq d}\) , with \(\eta \in K(C_{\leq d-1})\) , \(n_i \in Z\) and \(\dim(A/\mathfrak{p}_i) = d\) for \(1 \leqslant i \leqslant k\) ; one has \(\zeta_d(\xi) = \sum_{i=1}^{k} n_i[V(\mathfrak{p}_i)]\) and consequently \(\zeta_d(\xi) = 0\) implies \(\xi = i_d(\eta) \in \operatorname{Im}(i_d)\) , whence \(\operatorname{Ker}(\zeta_d) = \operatorname{Im}(i_d)\) .

One argues similarly for the second sequence.

Examples. --- 1) Suppose A is Noetherian and integral. Then we have \(Z^0 = \mathbf{Z}\){}{} ; one has \(C^{\geqslant 0} = C\) and \(z^0(M) = \operatorname{rg}(M)\) .[Spec(A)]. The modules of type \(C^{\geqslant 1}\) are therefore the torsion modules.

\begin{enumerate}
\def\labelenumi{\arabic{enumi})}
\setcounter{enumi}{1}
\item
  Suppose A is Noetherian and integrally closed. Then \(Z^1\) is identified with the group D(A) of divisors of A introduced in Chapter VII (§ 1, no. 3, th. 2, and no. 6, th. 3). The modules of type \(C^{\geqslant 2}\) are the pseudo-null modules (VII, § 4, no. 4, def. 2); if M is a finitely generated torsion module, then \(z^1(M) \in Z^1 = D(A)\) is the content \(\chi(M)\) of M (VII, § 4, no. 5, def. 4). Props. 10 and 11 of loc. cit. are therefore equivalent to the exactness of the sequence \(K(C^{\geqslant 2}) \to K(C^{\geqslant 1}) \to Z^1 \to 0\) .
\item
  The modules of type \(C_{\leqslant0}\) are the modules of dimension \(\leqslant0\) , that is, the modules of finite length (no. 4, remark 1). We have \(\text{long}_{\mathbf{A}}(\mathbf{M}) = \varepsilon(z_0(\mathbf{M}))\) for every A-module M of finite length, where \(\varepsilon: Z_0 \to Z\) associates to the linear combination \(\sum_{m} n_{m}[V(m)]\) the integer \(\sum_{m} n_{m}\) (IV, § 2, no. 5, corollary to prop. 8).
\item
  Suppose A is integral and of finite dimension. Set \(d = \dim(A)\) . Then one has \(\mathcal{C}_{\leq d} = \mathcal{C}\) , \(Z_d = \mathbf{Z}\){}{}\(Z^0\) , \(z_d(M) = \operatorname{rg}(M)\) . [Spec(A)] = \(z^0(M)\) , and the modules of type \(\mathcal{C}_{\leq d-1}\) are the torsion modules.
\end{enumerate}

